\documentclass[letterpaper,11pt]{article}

\usepackage[letterpaper,left=0.80in,right=0.80in,top=0.70in,bottom=0.70in]{geometry}
\usepackage[hidelinks]{hyperref}
\usepackage[english]{babel}
\input{glyphtounicode}

\pagestyle{empty}
\urlstyle{same}
\raggedright
\setlength{\parindent}{0pt}
\setlength{\parskip}{10pt}
\pdfgentounicode=1

\begin{document}

\begin{center}
  {\Large \textbf{Aryan Miriyala}}\\
  \small
  +1 419-315-0444 $|$
  \href{mailto:aryanmiriyala@gmail.com}{aryanmiriyala@gmail.com} $|$
  \href{https://www.linkedin.com/in/aryan-miriyala-7788a21b6/}{linkedin.com/in/aryan-miriyala} $|$
  \href{https://github.com/aryanmiriyala}{github.com/aryanmiriyala}
\end{center}

August 21, 2026

Google Hiring Team

Dear Hiring Team,

I am applying for the Software Engineer, Early Career, Campus role at Google. What stands out to me is the breadth of the role: building software that can move between web systems, information retrieval, distributed computing, AI, security, and large-scale user impact. My strongest matching work has been at the point where software has to be useful, testable, and dependable across real workflows.

At Actual Reality Technologies, I am improving a private Next.js and TypeScript customer portal with Firebase-backed application behavior, project-management API integration, authentication/session flows, dashboard status logic, and Vitest-tested customer/admin workflows. That work maps closely to Google's emphasis on designing, testing, deploying, maintaining, and improving software on small versatile teams. I have also used AI productivity tools including Codex and Claude Code in engineering workflows for planning, debugging, code review, and self-review while keeping human validation central.

My project background adds the systems and search side of the role. I built a C++/OpenMP adaptive runtime research system that uses telemetry and a UCB controller to tune scheduling policy, chunk size, and thread count across parallel workloads. I also built FalconGraph Search with a multithreaded C++ crawler, Python document processing, FAISS retrieval, FastAPI RAG endpoints, and a Next.js interface for cited answers. Those projects gave me hands-on experience with algorithms, Linux-based development, information retrieval, and software that has to stay understandable beyond a single feature.

I would welcome the opportunity to discuss how my experience across Python, C++, TypeScript, AI-assisted development, search systems, and secure full-stack workflows can contribute to Google's engineering teams.

Sincerely,

Aryan Miriyala

\end{document}
